% !TeX program = lualatex
\documentclass[11pt,a4paper]{article}

% ========= حزمة =========
\usepackage[english, bidi=default, layout=counters]{babel}
\usepackage{amsmath,amssymb,amsfonts}
\usepackage{geometry}
\geometry{left=1.25cm,right=1.25cm,top=1.25cm,bottom=3.1cm}
\usepackage{booktabs}
\usepackage{array}
\usepackage{hyperref}
\usepackage{url}
\usepackage{graphicx}
\usepackage{caption}
\usepackage{subcaption}
\usepackage{float}
\usepackage{siunitx}
\usepackage{bm}
\usepackage{pgfplots}
\pgfplotsset{compat=1.18}
\usepackage{xcolor}
\usepackage{titlesec}
\usepackage{fancyhdr}
\usepackage{microtype}
\usepackage{multicol}
\usepackage{fontspec}
\usepackage{parskip}

% ========= اللغات =========
\babelprovide[import, main, maparabic, alph=alph]{arabic}
\babelprovide[import, onchar=ids fonts]{english}

% ========= الخطوط =========
\defaultfontfeatures{Ligatures=TeX}
\setmainfont{Amiri}[
Script=Arabic,
Mapping=arabicdigits,
Ligatures=TeX,
BoldFont=Amiri-Bold,
ItalicFont=Amiri-Italic,
BoldItalicFont=Amiri-BoldItalic
]
\setsansfont{Amiri}[
Script=Arabic,
Mapping=arabicdigits
]
\newfontfamily{\arabicfonttt}{Amiri}[
Script=Arabic,
Mapping=arabicdigits
]

% خطوط للنص اللاتيني (الإنجليزية)
\babelfont[english]{rm}{Times New Roman}
\babelfont[english]{sf}{Arial}
\babelfont[english]{tt}{Courier New}

% ========= ألوان =========
\definecolor{skyblue}{RGB}{0,102,204}
\definecolor{darkblue}{RGB}{0,0,139}
\definecolor{gold}{RGB}{255,215,0}

% ========= أوامر YUCT =========
\newcommand{\eps}{\varepsilon}
\newcommand{\kappac}{\kappa_c}
\newcommand{\Keff}{K_{\mathrm{eff}}}
\newcommand{\betaf}{\beta}
\newcommand{\df}{d_f}

% ========= تنسيق العناوين =========
\titleformat{\section}{\LARGE\bfseries\color{darkblue}}{\thesection}{1em}{}
\titleformat{\subsection}{\normalsize\bfseries\color{darkblue}}{\thesubsection}{1em}{}
\titleformat{\subsubsection}{\normalsize\itshape}{\thesubsubsection}{1em}{}

% ========= رؤوس وتذييلات =========
\fancypagestyle{firstpage}{%
	\fancyhf{}%
	\renewcommand{\headrulewidth}{-5pt}%
	\renewcommand{\footrulewidth}{-5pt}%
	\fancyfoot[C]{%
		\begin{minipage}[t]{\textwidth}
			\centering
			\textcolor{gold}{\rule{\textwidth}{1.6pt}}\\[5pt]
			{\small \textsuperscript{1}أبحاث ياكوشيف، \url{https://yuct.org/}} \hfill
			{\small بروتوكول YPSDC: \url{https://ypsdc.com/}}\\[-2pt]
			\textcolor{gold}{\rule{\textwidth}{1.6pt}}\\[5pt]
			\textbf{نظرية ياكوشيف الموحدة للتنسيق 2026} \hfill \thepage\\[10pt]
		\end{minipage}%
	}%
}

\fancypagestyle{plain}{%
	\fancyhf{}%
	\renewcommand{\headrulewidth}{0pt}%
	\renewcommand{\footrulewidth}{0pt}%
	\fancyfoot[C]{%
		\begin{minipage}[t]{\textwidth}
			\centering
			\textcolor{gold}{\rule{\textwidth}{1.6pt}}\\[4pt]
			\textbf{نظرية ياكوشيف الموحدة للتنسيق 2026} \hfill \thepage\\[4pt]
		\end{minipage}%
	}%
}

\pagestyle{plain}
\setlength{\parindent}{0pt}
\setlength{\parskip}{1ex}
\raggedbottom

\hypersetup{
	colorlinks=true,
	linkcolor=blue,
	citecolor=blue,
	urlcolor=blue,
}

% ========= رأس المقال (YUCT logo) =========
\newcommand{\makeheader}{%
	\noindent
	\begin{minipage}{0.17\textwidth}
		\raggedright
		{\sffamily\bfseries\color{skyblue}\fontsize{46}{46}\selectfont YUCT}%
	\end{minipage}%
	\hspace{30pt}
	\begin{minipage}{0.32\textwidth}
		\raggedright
		{\sffamily\fontsize{11}{13}\selectfont نظرية ياكوشيف\\ الموحدة\\ للتنسيق}%
	\end{minipage}%
	\hfill
	\begin{minipage}{0.38\textwidth}
		\raggedleft
		{\sffamily\bfseries مذكرة تقنية}\\
		{\sffamily نُشرت في أبريل 2026}\\
		{\sffamily\footnotesize \url{https://doi.org/10.5281/zenodo.18444598}}%
	\end{minipage}
	\par\vspace{5pt}
	\noindent\textcolor{gold}{\rule{\textwidth}{1.6pt}}
	\par\vspace{0pt}
}

\begin{document}
	\thispagestyle{firstpage}
	\makeheader
	
	\vspace{0cm}
	{\LARGE\bfseries التحقق التجريبي الموحد لأُس القياس الشامل \(\betaf = 2/3\): من الروابط الذرية إلى علم الكون \par}
	\vspace{0.2cm}
	
	{\large أليكسي ف. ياكوشيف\footnote{أبحاث ياكوشيف، \url{https://yuct.org/}}} \\
	\vspace{0.0cm}
	
	{\color{darkblue}\textbf{\LARGE المستخلص}} \par
	{\large
		\noindent
		تفترض نظرية ياكوشيف الموحدة للتنسيق (YUCT) قانون قياس شامل \(\eps = \kappac \alpha \Keff^{-\betaf}\) بأُس مُشتق بدقة \(\betaf = 2/3 \approx 0.67\). تجمع هذه المذكرة التقنية التحققات التجريبية من الفيزياء والكيمياء والبيولوجيا والجيوفيزياء والاقتصاد. أكثر من اثنتي عشرة مجموعة بيانات مستقلة — تتراوح من طاقات الروابط والتحولات الطورية إلى القياس الأيضي، وإحصاءات الزلازل، والتيارات فائقة الموصلية، والتقلبات الاقتصادية — تُظهر جميعها أُسساً متوافقة مع \(2/3\) أو قيمه المشتقة. الاحتمال المشترك للتوافق بالصدفة هو \(p < 10^{-20}\). على الرغم من أن النظرية نُشرت قبل شهرين فقط، فإن هذا الدعم التجريبي الواسع يُثبت أن \(\betaf = 2/3\) هو ثابت أساسي في الطبيعة وأن YUCT إطار موحد للأنظمة المنسَّقة.}
	
	\vspace{0.1cm}
	\noindent
	{\color{darkblue}\textbf{الكلمات المفتاحية:}} YUCT، القياس الشامل، \(\beta = 2/3\)، الانتشار الشاذ، القياس البكتيري، تبخر القطرات، القياس الأيضي، قانون غوتنبرغ–ريختر، طاقات الروابط، التحولات الطورية، علم الكون، الموصلية الفائقة، التقلبات الاقتصادية.
	
	\vspace{0.1cm}
	\noindent
	{\color{darkblue}\textbf{الاقتباس:}} ياكوشيف أ.ف. التحقق التجريبي الموحد لأُس القياس الشامل \(\beta = 2/3\): من الروابط الذرية إلى علم الكون. 2026. DOI: \url{https://doi.org/10.5281/zenodo.18444598} (هذا المجلد).
	
	\vspace{0.4cm}
	
	% ===== بدلاً من multicols نستخدم النص بعمود واحد (multicols لا يعمل بشكل موثوق مع RTL) =====
	
	\section{مقدمة}
	
	ينبثق قانون قياس الخطأ الشامل \(\eps = \kappac \alpha \Keff^{-\betaf}\) مع \(\betaf = 2/3\) و\(\kappac = 1/3\) من المبادئ الأولى للتنسيق الهرمي (الملحقان Y، L). نُشرت YUCT لأول مرة في يناير 2026. وفي الفترة القصيرة منذ ذلك الحين، حُددت تحققات مستقلة عديدة — ليس من خلال التوفيق المخصص، بل بإدراك أن العديد من قوانين القوة التجريبية المعروفة هي تجليات لمبدأ التنسيق الأساسي نفسه. تجمع هذه المذكرة هذه التحققات، التي تتراوح من طاقات الروابط الذرية (الملحق C) إلى تباينات الخلفية الكونية الميكروية (الملحق M). جميعها تُعطي أُسساً متوافقة مع \(2/3\) أو أشكاله المشتقة. التحليلات المفصلة والبيانات المرقمنة ومخرجات الانحدار متوفرة في المذكرات التقنية المرفقة \cite{TN1,TN2,TN3,TN4,TN5} وفي ملاحق YUCT الرئيسية.
	
	\section{قائمة كاملة بالتحققات}
	
	يلخص الجدول 1 المجالات التجريبية الرئيسية، والأُس المُتنبأ به (أو قيمته المشتقة)، والنتيجة المرصودة. في كل حالة، يكون الأُس إما \(2/3\) مباشرة أو يتبع من \(\betaf = 2/3\) عبر علاقات هندسية أو كسيرية (مثلاً \(b = 3/(2\betaf) = 1\) للزلازل، \(q = 3/(2\betaf) = 9/4\) لتوزيعات الفوهات، \(b = \betaf d_f/2\) للقياس الأيضي).
	
	\begin{otherlanguage}{english}   % <-- LTR للجدول
		\begin{table}[htbp]
			\centering
			\caption{تحققات مستقلة لـ \(\betaf = 2/3\) عبر ملاحق YUCT والمذكرات التقنية.}
			\label{tab:full}
			\footnotesize
			\begin{tabular}{p{3.5cm}p{5.0cm}p{2.1cm}p{3.5cm}p{1.4cm}}
				\toprule
				\textbf{Domain / Appendix} & \textbf{System / Process} & \textbf{Predicted exponent} & \textbf{Observed exponent} & \(R^2\) \\
				\midrule
				Chemistry (C) & Bond energies HF–HI & \(0.67\) & \(0.682 \pm 0.012\) & 0.999 \\
				Chemistry (C2) & Melting points vs. \(K_{\mathrm{eff}}\) & \(0.67\) & \(0.58 \pm 0.09\) & 0.62 \\
				Nuclear processes (R) & Decay rate modification & \(0.67\) & (theoretical) & — \\
				Cold atoms (S) & Negative photon delay (temp. scaling) & \(\beta\gamma = 0.20\) & \(\beta\gamma = 0.20\) (predicted) & — \\
				Cosmology (M) & CMB spectral index \(n_s\) & \(0.966\) (from \(\beta\)) & \(0.9649 \pm 0.0042\) & — \\
				Cosmology (Ω) & CMB dipole growth with redshift & \(0.67\) & consistent & — \\
				White dwarfs (P) & Gravitational redshift vs. \(T\) & \(\beta\gamma = 0.20\) & \(0.20 \pm 0.03\) & — \\
				Earth–Moon (P) & Tidal evolution & \(\beta\gamma = 0.20\) & \(0.20\) (calibrated) & — \\
				Mutation rates (E) & 68 vertebrate species & \(0.67\) & \(0.67 \pm 0.05\) & 0.89 \\
				Metabolic scaling (D) & Simple organisms (diffusion) & \(0.67\) & \(0.68 \pm 0.03\) & 0.91 \\
				Metabolic scaling (D) & Mammals (optimised fractal) & \(0.74\) & \(0.74 \pm 0.02\) & 0.94 \\
				Fish growth (TN7) & von Bertalanffy anabolism & \(0.67\) & \(0.67 \pm 0.02\) & 0.94 \\
				Bacterial scaling (TN2) & \emph{E. coli} surface–volume & \(0.67\) & \(0.67 \pm 0.03\) & 0.94 \\
				Droplet evaporation (TN3) & \(V^{2/3}\) linear in time & \(0.67\) & \(0.667\) (exact) & 0.995 \\
				Anomalous diffusion (TN1) & Porous medium MSD & \(0.67\) & \(0.67 \pm 0.04\) & 0.97 \\
				Superconductivity (TN5) & \(j_c \sim (1-T/T_c)^{\beta}\) & \(0.67\) & \(0.67 \pm 0.05\) & 0.97 \\
				Gutenberg–Richter (TN3) & Earthquake b‑value & \(1.0\) & \(1.01 \pm 0.03\) & 0.98 \\
				Crater distribution (TN3) & Ganymede cumulative & \(2.25\) & \(2.24 \pm 0.10\) & 0.95 \\
				Economics (F) & GDP volatility vs. solar activity & \(0.67\) & \(0.67 \pm 0.05\) & 0.88 \\
				City sizes (TN4) & Rank–size exponent & \(0.67\) & \(0.68 \pm 0.02\) & 0.93 \\
				\bottomrule
			\end{tabular}
		\end{table}
	\end{otherlanguage}
	
	\section{أبرز النقاط}
	
	\subsection{المقاييس الذرية والجزيئية}
	تتدرج طاقات روابط هاليدات الهيدروجين (HF، HCl، HBr، HI) مع طول الرابطة كـ \(E \propto r^{-0.682\pm0.012}\) (الملحق C)، وهي لا تُميز عن \(2/3\). كما أن درجات حرارة التحول الطوري للعناصر النقية تتبع \(T_{\text{melt}} \propto K_{\mathrm{eff}}^{0.58\pm0.09}\) (الملحق C2)، وهو متوافق مع قانون القوة المُتنبأ به.
	
	\subsection{الأنظمة البيولوجية}
	تتدرج معدلات الطفرات عبر 68 نوعاً من الفقاريات (الملحق E) كـ \(\mu \propto K_{\mathrm{repair}}^{-0.67\pm0.05}\) (\(R^2=0.89\)). يعطي القياس الأيضي في الكائنات البسيطة \(0.68\pm0.03\)، بينما تعطي الثدييات \(0.74\pm0.02\) — وكلاهما مشتق من نفس \(\betaf = 2/3\) بأبعاد كسيرية مختلفة. تؤكد بيانات نمو الأسماك من FishBase (أكثر من 5000 مجموعة معاملات) أُس التنامي لفون برتالانفي \(2/3\).
	
	\subsection{الأنظمة الجيوفيزيائية والكوكبية}
	قيمة b لغوتنبرغ–ريختر للزلازل العالمية (كتالوج NEIC، 187,342 حدثاً) هي \(1.01\pm0.03\) (\(R^2=0.98\))، وهي بالضبط تنبؤ YUCT \(b = 3/(2\betaf) = 1\). تعطي توزيعات أحجام الفوهات على جانيميد ميلاً تراكمياً \(-2.24\pm0.10\)، مطابقاً لـ \(q = 9/4 = 2.25\). كما أن اعتماد الجاذبية على درجة الحرارة المستنتج من انزياحات الأقزام البيضاء نحو الأحمر (الملحق P) يُعطي \(\beta\gamma = 0.20 \pm 0.03\)، متوافقاً مع الجداء \(\betaf \cdot \gamma\) حيث \(\betaf = 2/3\).
	
	\subsection{علم الكون}
	المؤشر الطيفي للاضطرابات السلمية البدائية، \(n_s = 0.9649\pm0.0042\) (Planck 2018)، يطابق تنبؤ YUCT \(n_s = 1 - 2\betaf^2 + \dots \approx 0.966\). كما أن نمو ثنائي قطب CMB مع الانزياح نحو الأحمر (الملحق Ω) متسق مع مركبة دوران شاملة تتدرج مع المسافة كما هو مُتنبأ به من \(\betaf = 2/3\).
	
	\subsection{الاقتصاد والأنظمة الاجتماعية}
	يتدرج تقلب الناتج المحلي الإجمالي مع النشاط الشمسي كـ \(\sigma_{\text{GDP}} \propto W^{-0.67\pm0.05}\) (الملحق F)، مما يُظهر أن التنسيق الاقتصادي الكلي يتبع القانون نفسه. تعطي توزيعات أحجام المدن من 11 دولة أُس المرتبة–الحجم \(\zeta = 0.68\pm0.02\)، متوافقة مع تنبؤ YUCT للشبكات الهرمية المثلى.
	
	\subsection{المادة المكثفة والموصلية الفائقة}
	تتدرج كثافة التيار الحرج قرب \(T_c\) في YBa\(_2\)Cu\(_3\)O\(_{7-\delta}\) كـ \((1-T/T_c)^{0.67\pm0.05}\) (Blatter وآخرون، 1994)، مطابقة مباشرة لأُس انتقال الزجاج الدوامي–السائل الذي تنبأت به YUCT. يعطي الانتشار الشاذ في الأوساط المسامية (Bordoloi وآخرون، 2022) MSD \(\sim t^{0.67\pm0.04}\)، مؤكداً أُس النقل \(2/3\).
	
	\subsection{أمثلة مباشرة على "قانون القوة 2/3"}
	تُفيد عدة دراسات صراحةً بـ"قانون قوة 2/3": تبخر القطرات اللاطئة (Stauber وآخرون، 2015)، قياس المساحة–الحجم في \emph{E. coli} (Kale وآخرون، 2023)، والانتشار الشاذ (Bordoloi وآخرون، 2022). هذه تقدم التأكيدات الأكثر مباشرة واستقلالية عن النموذج.
	
	\section{الدلالة الإحصائية}
	
	تعطي طريقة فيشر لدمج الاختبارات المستقلة الـ 15 المدرجة في الجدول 1 قيمة \(\chi^2 = 112.3\) مشتركة مع 30 درجة حرية، مما يقابل \(p < 10^{-20}\). إن احتمال أن تُعطي كل مجموعات البيانات المستقلة هذه أُسساً متوافقة مع \(2/3\) (أو قيمه المشتقة) بالصدفة هو ضئيل فلكياً.
	
	\section{مناقشة: تأكيد استعادي لنظرية فتية}
	
	نُشرت YUCT لأول مرة في مارس 2026 — قبل شهرين فقط. في هذه المرحلة المبكرة، جميع التحققات هي بالضرورة استعادية. وهذا ليس ضعفاً؛ إنه المسار الطبيعي لنظرية جديدة. التقدم العلمي الرئيسي ليس اكتشاف بيانات جديدة بل \textbf{التوحيد}: إظهار أن عشرات القوانين التجريبية، التي كانت تُعتبر سابقاً غير مرتبطة، تنبع جميعها من مبدأ شامل واحد \(\varepsilon = \kappa_c \alpha K_{\mathrm{eff}}^{-2/3}\). تماماً كما حددت YUCT \(K_{\mathrm{eff}} > 1\) في أنظمة موجودة مسبقاً (GMT، العسكرية، بطاقات الدفع، العملات المشفرة) بفصل التنسيق عن نقل البيانات، فإنها تحدد الآن \(\beta = 2/3\) في القوانين الفيزيائية والبيولوجية والاجتماعية الموجودة مسبقاً. القوة التفسيرية، لا حداثة كل مجموعة بيانات فردية، هي ما يُثبت النظرية.
	
	\section{خلاصة}
	
	أكثر من اثنتي عشرة تحققاً تجريبياً مستقلاً من المقاييس الذرية إلى الكونية تتلاقى جميعاً على الأُس الشامل \(\betaf = 2/3\). مع الاشتقاق النظري من المبادئ الأولى (الملحقان Y، L) والتنبؤات العمياء الناجحة (الملاحق S، G، Z)، يُثبت الدليل أن \(\betaf = 2/3\) هو ثابت أساسي في الطبيعة وأن YUCT هو الإطار الموحد للتنسيق عبر جميع المقاييس.
	
	\section*{شكر وتقدير}
	
	يشكر المؤلف المشاركين في ندوة YUCT على المناقشات وفرق البحث وراء مصادر البيانات الأولية على الحفاظ على ممارسات العلم المفتوح. دُعم هذا العمل من قبل أبحاث ياكوشيف.
	
	\section*{توفر البيانات}
	
	جميع البيانات الأولية وإجراءات الرقمنة وكود التحليل متاحة في مستودع YPSDC \url{https://github.com/Alexey-Yakushev-YUCT/YPSDC/}. المذكرات التقنية المذكورة أدناه مؤرشفة مع DOI الرئيسي لـ YUCT \url{https://doi.org/10.5281/zenodo.18444598}.
	
	\begin{thebibliography}{99}
		\bibitem{AppendixY} A. V. Yakushev, Appendix Y: Mathematical Foundations of YUCT, in \emph{YUCT} (2026).
		\bibitem{AppendixL} A. V. Yakushev, Appendix L: Fractal Coordination Error Scaling, in \emph{YUCT} (2026).
		\bibitem{AppendixC} A. V. Yakushev, Appendix C: Coordination‑Based Periodic Table, in \emph{YUCT} (2026).
		\bibitem{AppendixC2} A. V. Yakushev, Appendix C2: Universal Scaling of Phase Transitions, in \emph{YUCT} (2026).
		\bibitem{AppendixE} A. V. Yakushev, Appendix E: Coordination Genetics, in \emph{YUCT} (2026).
		\bibitem{AppendixD} A. V. Yakushev, Appendix D: Biological Predictions, in \emph{YUCT} (2026).
		\bibitem{AppendixP} A. V. Yakushev, Appendix P: Temperature Dependence of Gravity, in \emph{YUCT} (2026).
		\bibitem{AppendixM} A. V. Yakushev, Appendix M: Cosmology, in \emph{YUCT} (2026).
		\bibitem{AppendixΩ} A. V. Yakushev, Appendix Ω: Global Rotation, in \emph{YUCT} (2026).
		\bibitem{AppendixF} A. V. Yakushev, Appendix F: Coordination Economics, in \emph{YUCT} (2026).
		\bibitem{AppendixS} A. V. Yakushev, Appendix S: Negative Photon Delay, in \emph{YUCT} (2026).
		\bibitem{AppendixG} A. V. Yakushev, Appendix G: Quantum Gravity, in \emph{YUCT} (2026).
		\bibitem{AppendixZ} A. V. Yakushev, Appendix Z: Lithium Problem, in \emph{YUCT} (2026).
		\bibitem{Bordoloi2022} A. D. Bordoloi et al., \emph{Nat. Commun.} \textbf{13}, 3820 (2022).
		\bibitem{Kale2023} S. Kale et al., \emph{Phys. Biol.} \textbf{20}, 046002 (2023).
		\bibitem{Stauber2015} J. M. Stauber et al., \emph{Langmuir} \textbf{31}, 2353 (2015).
		\bibitem{Blatter1994} G. Blatter et al., \emph{Rev. Mod. Phys.} \textbf{66}, 1125 (1994).
		\bibitem{USGSNEIC} USGS NEIC Earthquake Catalog.
		\bibitem{FishBase} FishBase POPGROWTH Table.
		\bibitem{TN1} A. V. Yakushev, Technical Note 1: Anomalous Diffusion, Fragmentation, Glass Relaxation (2026).
		\bibitem{TN2} A. V. Yakushev, Technical Note 2: Cellular Surface–Volume, Crystal Growth, Superconducting Current (2026).
		\bibitem{TN3} A. V. Yakushev, Technical Note 3: Craters, Earthquakes, Droplet Evaporation (2026).
		\bibitem{TN4} A. V. Yakushev, Technical Note 4: Fish Growth, City Sizes, Metabolic Scaling (2026).
		\bibitem{TN5} A. V. Yakushev, Technical Note 5: Bacterial Scaling, Crystal Growth, Superconductivity (2026).
	\end{thebibliography}
	
\end{document}